\chapter{Introduction}

The \textbf{Pinacoteca} project delivers a management system for a gallery, covering access control, climate, humidity, lighting, and LCD visualization. The goal is to integrate sensors and actuators into a modular Arduino firmware, with unit tests runnable on a PC and a repeatable build flow.

This documentation is meant to describe both the technical solution and the rationale behind design choices. Each section dives into the behavior step by step, highlighting safety aspects, adopted trade-offs, and the reasoning behind the exposed APIs.

The main features are:
\begin{itemize}
  \item counting entrances and exits with a servo-assisted turnstile;
  \item indicating available capacity with a LED traffic light;
  \item temperature regulation with heating or cooling modes;
  \item humidity control with threshold and actuator (LED in simulation);
  \item ambient light control with a photoresistor and dimmer;
  \item a 16x2 LCD panel with informative pages.
\end{itemize}

The code is structured into independent modules under \texttt{lib/}. This choice reduces cross-dependencies, makes behavior easier to read, and simplifies validation through unit tests.

The system is built around reliability: if a sensor reading is invalid, control modules switch to a fallback behavior that turns outputs off, avoiding unsafe or non-deterministic states.
